\chapter{Background}
%%%%%%%%%%%%%%%%%%%%%%

To understand the design and implementation of an adaptive actuation system, it is necessary to establish the technical foundations of continuous sensing hardware, the mathematical logic of modern trigger parameters, and the biomechanical constraints of the human hand.

\section{Continuous Depth-Sensing Mechanisms}

Traditional mechanical switches operate on a binary principle: a circuit is either open or closed based on a physical contact point. In contrast, continuous depth-sensing switches measure the absolute position of the stem throughout its entire travel distance.

\subsection{Hall-Effect and Magnetic Sensing}
Hall-effect sensors utilize a permanent magnet attached to the switch stem and a sensor on the Printed Circuit Board (PCB). As the magnet approaches the sensor, the Hall voltage varies proportionally to the magnetic field strength~\cite{wooting_rapid_trigger, pcgamer_hall_effect}. This allows the firmware to map voltage levels to precise displacement values in millimeters. Compared to classic contact switches, this technology reduces dependence on debounce filtering and enables lower-latency signal processing.

\subsection{Haptic Inductive Trigger System (HITS)}
Refined in the 2026 Logitech SUPERSTRIKE series, the Haptic Inductive Trigger System (HITS) replaces magnets with an inductive coil architecture~\cite{logitech_hits_patent, techradar_superstrike_2026}. By measuring changes in electromagnetic induction caused by the proximity of a metallic puck within the trigger, HITS provides higher resolution and lower susceptibility to external magnetic interference compared to Hall-effect sensors. Furthermore, HITS integrates localized haptic motors to simulate tactile "bumps" or "clicks" at software-defined depths, decoupling the physical feel of the switch from its electrical actuation point~\cite{fox_superstrike_2026}.

\section{Control Parameters: AP and RT}

The transition to continuous sensing introduced two primary variables that define the "feel" and speed of an input device.

\subsection{Actuation Point (AP)}
The Actuation Point (AP) is the specific depth ($d$) at which a button-press event is registered. In traditional switches, this is fixed (e.g., 2.0mm). In analog systems, AP is a variable $d_{act} \in [0.1, 4.0]$mm. Lowering $d_{act}$ reduces mechanical pre-travel, allowing for faster response times, but increases the risk of accidental actuations caused by resting finger weight.

\subsection{Rapid Trigger (RT)}
Rapid Trigger (RT) eliminates the fixed reset point found in mechanical switches. In a standard switch, the button must travel back up past a specific reset point before it can be pressed again. RT logic instead monitors the direction of travel. As soon as the stem moves upward by a specific threshold ($\Delta_{up}$), the input resets. Conversely, as soon as it moves back down by $\Delta_{down}$, it reactuates~\cite{wooting_rapid_trigger, drone_rapid_trigger}. This allows rapid repeated clicking, which is critical in high-APM genres.

\section{Biomechanical Variance and Neuromotor Fatigue}

The efficiency of AP and RT is limited by the physiological state of the user. Human motor control is subject to "signal-dependent noise," where the variance of movement increases with the force or speed of the action.

\subsection{Resting Pressure and Tremors}
Even at rest, users exert varying degrees of force on a peripheral. Factors such as cognitive stress or caffeine intake can cause micro-tremors-high-frequency, low-amplitude oscillations in displacement telemetry~\cite{stress_pressure}. If the AP is set to 0.1mm and the user's tremor amplitude is 0.15mm, the system will register "ghost" inputs.

\subsection{Neuromotor Fatigue Signatures}
Extended sessions lead to fatigue in the \textit{flexor digitorum superficialis} and \textit{extensor digitorum} muscles~\cite{fatigue_flexor, clicking_fatigue_ergonomics_2024}. Fatigue manifests in two measurable telemetry changes:
\begin{enumerate}
    \item \textbf{Velocity Decay:} The "snap-back" velocity (the speed at which a finger releases a key) decreases as the extensor muscles tire~\cite{tholl_wrist_fatigue_2025}.
    \item \textbf{Resting Weight Shift:} As the flexor muscles lose tonicity, the resting weight of the finger often increases, leading to deeper "resting" displacement in the analog sensor~\cite{heimhofer2024finger, 1hp_esports_ergonomics}.
\end{enumerate}

\section{On-Device Inference and TinyML (Scope Context)}

This thesis does not implement or evaluate on-device inference, but brief context is included to position future deployment constraints. Continuous telemetry at high polling rates imposes tight latency budgets that are difficult to satisfy through host-side software alone because of USB polling and OS scheduling behavior~\cite{monaco_sok_2018}. Recent peripheral MCUs, including RISC-V-class devices, are increasingly capable of quantized inference~\cite{tinyml_mcu_inference, tinyml_edge_2025, singh2023edge}. In principle, INT8 models could support pre-report filtering under sub-millisecond deadlines; in practice, validating this claim requires firmware-level access and hardware-in-the-loop measurement, which are out of scope for the present study.

\section{Signal Processing and Optimisation Foundations}
\label{sec:bg-methods}

The adaptive calibration pipeline draws on several established signal processing and machine learning methods. This section describes their mathematical properties; their problem-specific justifications appear in the Design chapter.

\subsection{Savitzky-Golay Smoothing}
Savitzky-Golay (SG) filtering fits a polynomial of degree $p$ to each sliding window of $W$ samples in a least-squares sense and evaluates the polynomial at the window centre~\cite{savitzky_golay_1964}. Because the convolution coefficients are precomputed from the normal equations, the runtime cost is $O(n)$ regardless of window width, identical to a moving average. Crucially, the polynomial fit exactly reproduces moments up to degree $p$, so peak positions and edge inflection points in the signal are not distorted—a property absent from symmetric Gaussian or box filters, which both shift extrema toward the window centre.

\subsection{Hermite Cubic Interpolation: Akima and Makima}
Akima interpolation constructs a piecewise cubic Hermite spline whose tangent at each knot is a weighted average of the four adjacent finite differences, with weights inversely proportional to the absolute difference between consecutive slopes~\cite{akima_1970}. This weighting isolates the tangent estimate from outlier slopes: a single anomalous interval does not propagate ringing across the spline, unlike a natural cubic spline whose tangents are coupled globally. Akima does not enforce monotonicity, so it can represent non-monotone smooth curves. Modified Akima (Makima) adds a half-mean-slope regularisation term to the weight formula, preventing zero-weight degeneracy when two adjacent slopes are equal and opposite—a degenerate case that arises frequently in quantized sensor signals.

\subsection{Covariance Matrix Adaptation Evolution Strategy}
CMA-ES is a derivative-free evolutionary optimiser for continuous search spaces~\cite{hansen_cma_2016}. It maintains a multivariate Gaussian distribution over candidate solutions and, at each generation, moves the distribution mean toward the best-performing candidates and adapts the full covariance matrix to reflect the local curvature of the objective landscape. This covariance adaptation lets the algorithm exploit correlations between parameters and navigate elongated, ill-conditioned valleys that cause gradient-free methods with diagonal step sizes (e.g.\ Nelder-Mead or coordinate descent) to stall. CMA-ES converges in $O(n^2)$ function evaluations for an $n$-dimensional problem, far fewer than the $O(N^n)$ evaluations required by grid search over a width-$N$ grid.

\subsection{Ensemble Methods: Random Forests and Gradient Boosting}
A Random Forest~\cite{breiman_rf_2001} builds an ensemble of decision trees, each trained on a bootstrapped subsample with a random subset of features considered at each split, and averages their posterior class probabilities. The double randomisation decorrelates tree errors, yielding ensemble variance lower than any individual tree without significantly increasing bias. Gradient Boosted Trees~\cite{friedman_gb_2001} instead fit trees sequentially: each tree is trained to correct the pseudo-residual of the current ensemble via gradient descent on a differentiable loss. This staged correction reduces bias relative to a Random Forest, at the cost of longer training time and greater sensitivity to hyperparameter tuning.

\subsection{Bayesian Optimisation and Tree-Structured Parzen Estimators}
Bayesian optimisation builds a probabilistic surrogate model of the objective function and selects the next evaluation point by maximising an acquisition function that balances exploration of uncertain regions against exploitation of known good regions. The Tree-structured Parzen Estimator (TPE)~\cite{bergstra_tpe_2011} implements this by modelling two separate densities over the hyperparameter space: $\ell(x)$ over configurations that produced good results, and $g(x)$ over those that produced poor results. Candidates are drawn from $\ell(x)$ and ranked by the ratio $\ell(x)/g(x)$, concentrating search in regions that distinguish good from bad performance. The Optuna framework~\cite{optuna_akiba_2019} implements TPE with asynchronous evaluation and persistent study storage backed by SQLite.


%%%%%%%%%%%%%%%%
